% Part: computability
% Chapter: computability-theory
% Section: fixed-point-thm

\documentclass[../../../include/open-logic-section]{subfiles}

\begin{document}

% SEGMENT OLP-0249-S01

\olfileid{cmp}{thy}{fix}
\olsection{நிலைப்புள்ளித் தேற்றம்}

நிற்றல் சிக்கலை மீண்டும் கருதுவோம். தற்காலிகக் குறியீடாக,
$\tuple{x, y}$ என்பதற்கு $\gn{\cfind{x}(y)}$ என்று எழுதுவோம்;
இதை $\cfind{x}(y)$ என்ற மதிப்புக்கான ஒரு ``பெயரைக்'' குறிப்பதாகக்
கருதுக. இக்குறியீட்டைக் கொண்டு நிற்றல் சிக்கல் தீர்மானிக்கவியலாதது
என்பதற்கான நமது நிரூபணங்களில் ஒன்றை வேறுவிதமாகக் கூறலாம்.

கேள்வி: ஒவ்வொரு $x$, $y$ க்கும் பின்வரும் பண்பைக் கொண்ட
கணிக்கத்தக்க சார்பு $h$ ஒன்று உள்ளதா?
\[
h(\gn{\cfind{x}(y)}) =
\begin{cases}
1 & \text{$\cfind{x}(y) \fdefined$ எனில்} \\
0 & \text{இல்லையெனில்.}
\end{cases}
\]

விடை: இல்லை; இருந்தால், பகுதிச் சார்பு
\[
g(x) \simeq
\begin{cases}
0 & \text{$h(\gn{\cfind{x}(x)}) = 0$ எனில்} \\
\text{வரையறுக்கப்படவில்லை} & \text{இல்லையெனில்}
\end{cases}
\]
கணிக்கத்தக்கதாக இருக்கும்; எனவே அதற்கு ஏதோ சுட்டெண்~$e$ இருக்கும்.
ஆனால் அப்போது
\[
\cfind{e}(e) \simeq
\begin{cases}
0 & \text{$h(\gn{\cfind{e}(e)}) = 0$ எனில்} \\
\text{வரையறுக்கப்படவில்லை} & \text{இல்லையெனில்,}
\end{cases}
\]
என்பது கிடைக்கும். இந்நேர்வில் $\cfind{e}(e)$ வரையறுக்கப்பட்டால்,
எனில் மட்டுமே, அது வரையறுக்கப்படவில்லை; இது முரண்பாடு.

இப்போது $\cfind{e}$ உள்ள சமன்பாட்டைப் பார்க்க. ஓர் அர்த்தத்தில்
அங்கு தன்குறிப்பு உள்ளது: $\cfind{e}(e)$ இன் மதிப்பு ஒரு குறிப்பிட்ட
முறையில் $\gn{\cfind{e}(e)}$ ஐச் சாருமாறு அமைத்துள்ளோம்.
முரண்பாடுகளை நிறுவுவதற்கு மட்டுமல்லாமல், பொதுவாகவே இதை நாம்
{\em செய்ய முடியும்} என்று நிலைப்புள்ளித் தேற்றம் கூறுகிறது.

நிலைப்புள்ளித் தேற்றத்தைக் கூறுவதற்கான சமானமான இரு வழிகளை
\olref{lem:fixed-equiv} வழங்குகிறது. தருக்கவியல் நோக்கில், இரு கூற்றுகளும்
உண்மையாக இருப்பதிலிருந்து அவை சமானமானவை என்பது பின்வருகிறது.
ஆனால் உண்மையில் நாம் கூறுவது, ஒவ்வொன்றிலிருந்தும் மற்றொன்று நேரடியாகப்
பின்வருவதால், அவற்றை ஒரே தேற்றத்தின் மாற்றுக் கூற்றுகளாக எடுத்துக்கொள்ளலாம்
என்பதே.

\begin{lem}
\ollabel{lem:fixed-equiv}
பின்வரும் கூற்றுகள் சமானமானவை:
\begin{enumerate}
\item ஒவ்வொரு பகுதி கணிக்கத்தக்க சார்பு $g(x,y)$ க்கும்,
  ஒவ்வொரு~$y$ க்கும் பின்வருமாறு அமையும் சுட்டெண்~$e$ ஒன்று உள்ளது:
\[
\cfind{e}(y) \simeq g(e,y).
\]
\item ஒவ்வொரு கணிக்கத்தக்க சார்பு~$f(x)$ க்கும்,
  ஒவ்வொரு~$y$ க்கும் பின்வருமாறு அமையும் சுட்டெண்~$e$ ஒன்று உள்ளது:
\[
\cfind{e}(y) \simeq \cfind{f(e)}(y).
\]
\end{enumerate}
\end{lem}

\begin{proof}
$(1) \Rightarrow (2)$: $f$ கொடுக்கப்பட்டால்,
$g(x,y) \simeq \fn{Un}(f(x),y)$ என்று $g$ ஐ வரையறுக்க.
ஒவ்வொரு~$y$ க்கும் பின்வருமாறு அமையும் சுட்டெண்~$e$ ஐப் பெற
(1) ஐப் பயன்படுத்துக:
\begin{align*}
\cfind{e}(y) & = \fn{Un}(f(e),y) \\
& = \cfind{f(e)}(y).
\end{align*}

$(2) \Rightarrow (1)$: $g$ கொடுக்கப்பட்டால், ஒவ்வொரு $x$,~$y$
க்கும் $\cfind{f(x)}(y) \simeq g(x,y)$ என அமையும் $f$ ஐப் பெற
$s$-$m$-$n$ தேற்றத்தைப் பயன்படுத்துக. பின்வருமாறு அமையும்
சுட்டெண்~$e$ ஐப் பெற (2) ஐப் பயன்படுத்துக:
\begin{align*}
\cfind{e}(y) & = \cfind{f(e)}(y) \\
& = g(e,y).
\end{align*}
இத்துடன் நிரூபணம் முடிகிறது.
\end{proof}

\begin{explain}
கூற்று (1) உண்மை (எனவே (2) உம் உண்மை) என்று காட்டுவதற்கு முன்,
அது எவ்வளவு விந்தையானது என்பதைக் கருதுக. $e$ ஐ ஒரு கணினி நிரலாகக்
கருதுக. எந்தப் பகுதி கணிக்கத்தக்க $g(x,y)$ கொடுக்கப்பட்டாலும்,
$g_e(y) \simeq g(e,y)$ ஐக் கணிக்கும் கணினி நிரல் $e$ ஒன்றை
கண்டுபிடிக்கலாம் என்று கூற்று (1) சொல்கிறது. வேறுவிதமாகக் கூறினால்,
தன்னையே குறிப்பிட்டு ஒரு சார்பைக் கணிக்கும் கணினி நிரலைக்
கண்டுபிடிக்கலாம்.
\end{explain}

\begin{thm}
\olref{lem:fixed-equiv} இலுள்ள இரு கூற்றுகளும் உண்மையானவை.
குறிப்பாக, ஒவ்வொரு பகுதி கணிக்கத்தக்க சார்பு $g(x,y)$ க்கும்,
ஒவ்வொரு~$y$ க்கும் பின்வருமாறு அமையும் சுட்டெண்~$e$ ஒன்று உள்ளது:
\[
\cfind{e}(y) \simeq g(e,y).
\]
\end{thm}

\begin{proof}
தேவையான கூறுகள் மேலுள்ள நிற்றல் சிக்கல் பற்றிய விவாதத்தில் ஏற்கெனவே
மறைமுகமாக உள்ளன. ஒவ்வொரு $x$ க்கும்
$f_x(y) \simeq \cfind{x}(x,y)$ என்ற சார்புக்கான ஒரு சுட்டெண்ணைத்
திருப்பித்தரும் கணிக்கத்தக்க சார்பாக $\fn{diag}(x)$ ஐ எடுத்துக்கொள்க;
அதாவது
\[
\cfind{\fn{diag}(x)}(y) \simeq \cfind{x}(x,y).
\]
2-இடச் சார்புக்கான நிரலை, அந்நிரலையே அதன் முதல் மதிப்புருவாக
நிலைநிறுத்திப் பெறப்படும் 1-இடச் சார்புக்கான நிரலாக மாற்றும்
சார்பாக $\fn{diag}$ ஐக் கருதுக. $\fn{diag}$ சார்பை முறைசார்ந்து
பின்வருமாறு வரையறுக்கலாம்: முதலில்
\[
s(x,y) \simeq \fn{Un}^2(x,x,y),
\]
என்று $s$ ஐ வரையறுக்க; இங்கு $\fn{Un}^2$ என்பது பகுதி கணிக்கத்தக்க
2-இடச் சார்புகளுக்கான அனைத்தளாவிய 3-இடச் சார்பு. பிறகு
$s$-$m$-$n$ தேற்றத்தால்,
\[
\cfind{\fn{diag}(x)}(y) \simeq s(x,y).
\]
என்பதை நிறைவேற்றும் முதன்மை மீள்வரையறைச் சார்பு $\fn{diag}$
ஒன்றைக் கண்டுபிடிக்கலாம்.

இப்போது சார்பு $l$ ஐ
\[
l(x,y) \simeq g(\fn{diag}(x),y).
\]
என்று வரையறுத்து, $l$ கான ஒரு சுட்டெண்ணாக $\gn{l}$ ஐ எடுத்துக்கொள்க.
இறுதியாக $e = \fn{diag}(\gn{l})$ என்க. அப்போது ஒவ்வொரு $y$ க்கும்,
\begin{align*}
\cfind{e}(y) & \simeq \cfind{\fn{diag}(\gn{l})}(y) \\
& \simeq \cfind{\gn{l}}(\gn{l}, y) \\
& \simeq l(\gn{l}, y) \\
& \simeq g(\fn{diag}(\gn{l}),y) \\
& \simeq g(e, y),
\end{align*}
என்று வேண்டியவாறு கிடைக்கிறது.
\end{proof}

\begin{explain}
இங்கு என்ன நடக்கிறது? தன்னையே அச்சிடும் ஒரு கணினி நிரலை எழுதும்
பணி உங்களுக்குக் கொடுக்கப்பட்டுள்ளது என்க. மேலும், வளமானதும்
விந்தையானதுமான எழுத்துத்தொடர்ச் சார்புகளின் நூலகம் கொண்ட ஒரு
நிரலாக்க மொழியில் நீங்கள் பணியாற்றுவதாகக் கொள்க. குறிப்பாக,
உங்கள் நிரலாக்க மொழியில் பின்வருமாறு செயல்படும் சார்பு
$\fn{diag}$ ஒன்று இருப்பதாகக் கொள்க: உள்ளீட்டு எழுத்துத்தொடர்~$s$
கொடுக்கப்பட்டால், ~$s$ இல் தோன்றும் `x' குறியின் ஒவ்வொரு நிகழ்வையும்
$\fn{diag}$ கண்டறிந்து, மூல எழுத்துத்தொடரின் மேற்கோளிடப்பட்ட
வடிவத்தால் அதை மாற்றுகிறது. எடுத்துக்காட்டாக,
\begin{quote}
\begin{verbatim}
hello x world
\end{verbatim}
\end{quote}
என்ற எழுத்துத்தொடரை உள்ளீடாகக் கொடுத்தால், சார்பு
\begin{quote}
\begin{verbatim}
hello 'hello x world' world
\end{verbatim}
\end{quote}
என்பதை வெளியீடாகத் தருகிறது. அப்போது வேண்டிய நிரலை எழுதுவது எளிது;
\begin{quote}
\begin{verbatim}
print(diag('print(diag(x))'))
\end{verbatim}
\end{quote}
என்பது வேண்டியதைச் செய்கிறது என்று சரிபார்க்கலாம். C++, Java போன்ற
பொதுவான நிரலாக்க மொழிகளிலும் இதே கருத்து (மேலும் விரிவான
நிறைவேற்றத்துடன்) செயல்படுகிறது.

நிலைப்புள்ளித் தேற்றத்தின் நிரூபணத்திலிருந்து இன்னும் ஓரிரு படிகளே
தொலைவில் இருக்கிறோம். அச்சிடும் சார்பின் ஒரு மாற்றுரு
$\fn{print}(x,y)$, எழுத்துத்தொடர் $x$ ஐயும் மற்றொரு எண்சார்
மதிப்புரு $y$ ஐயும் ஏற்று, $x$ என்ற எழுத்துத்தொடரை $y$ முறை
அச்சிடுவதாகக் கொள்க. அப்போது ``நிரல்''
\begin{quote}
\begin{verbatim}
getinput(y);
print(diag('getinput(y); print(diag(x), y)'), y)
\end{verbatim}
\end{quote}
உள்ளீடு $y$ இல் தன்னையே $y$ முறை அச்சிடுகிறது.
$\fn{getinput}$---$\fn{print}$---$\fn{diag}$ என்ற எலும்புக்கூட்டை
ஏதேனும் சார்பு $g(x,y)$ ஆல் மாற்றினால்
\begin{quote}
\begin{verbatim}
g(diag('g(diag(x), y)'), y)
\end{verbatim}
\end{quote}
என்பது கிடைக்கிறது. இது உள்ளீடு $y$ இல், நிரல் தன்னையும் $y$ ஐயும்
உள்ளீடுகளாகக் கொண்டு $g$ ஐ இயக்கும் நிரல். ``மேற்கோளிடுதல்''
என்பதை ``ஒரு சுட்டெண்ணைப் பயன்படுத்துதல்'' என்று கருதினால்,
மேலுள்ள நிரூபணம் கிடைக்கிறது.

இப்போதைக்கு, இந்நிரூபணத்தை முறைசார்ந்த தந்திரம் அல்லது கருமந்திரம்
என்று நீங்கள் கருத விரும்பினால் அது சரிதான். ஆனால் மேலே கொடுக்கப்பட்ட
வாதத்தின் விவரங்களை உங்களால் மீளமைக்க முடிய வேண்டும். முழுமையின்மைத்
தேற்றங்களையும் (தொடர்புடைய ``நிலைப்புள்ளித் தேற்றத்தையும்'')
நிறுவும்போது, இது ஏன் செயல்படுகிறது என்பதைப் புரிந்துகொள்வதற்கான
வேறு வழிகளை விவாதிப்போம்.
\end{explain}

\begin{tagblock}{lambda}
\begin{digress}
இதே கருத்தைப் பயன்படுத்தி ஒரு ``நிலைப்புள்ளிச்'' சேர்ப்பானைப்
பெறலாம். லாம்டா உறுப்பு $g$ ஒன்று உங்களிடம் உள்ளது என்றும்,
$k$ ஆனது $gk$ க்கு $\beta$-சமானம் என்ற பண்பைக் கொண்ட மற்றோர்
உறுப்பு $k$ வேண்டும் என்றும் கொள்க. நமது குறியீட்டு மரபுகளைப்
பயன்படுத்தி
\[
\fn{diag}(x) = xx
\]
மற்றும்
\[
l(x) = g(\fn{diag}(x))
\]
என்ற உறுப்புகளை வரையறுக்க; வேறுவிதமாகக் கூறினால், $l$ என்பது
$\lambd[x][g(xx)]$ என்ற உறுப்பு. $k$ ஆனது $ll$ என்ற உறுப்பாக
இருக்கட்டும். அப்போது
\begin{align*}
k & = (\lambd[x][g(xx)])(\lambd[x][g(xx)]) \\
& \red  g((\lambd[x][g(xx)])(\lambd[x][g(xx)])) \\
& = gk.
\end{align*}
பின்வருமாறு எடுத்துக்கொண்டால்
\[
Y = \lambd[g][((\lambd[x][g(xx)])(\lambd[x][g(xx)]))]
\]
$Yg$ உம் $g(Yg)$ உம் ஒரு பொதுவான உறுப்புக்குக் குறைகின்றன;
எனவே $Yg \equiv_\beta g(Yg)$. இது ``கரியின் சேர்ப்பான்''
என்று அழைக்கப்படுகிறது. இதற்குப் பதிலாக
\[
Y = (\lambd[xg][g(xxg)])(\lambd[xg][g(xxg)])
\]
என்று எடுத்துக்கொண்டால், $Yg$ ஆனது உண்மையில் $g(Yg)$ குக் குறைகிறது;
இது மேலும் வலுவான கூற்று. $Y$ இன் இந்தப் பிந்தைய வடிவம்
``டியூரிங்கின் சேர்ப்பான்'' என்று அழைக்கப்படுகிறது.
\end{digress}
\end{tagblock}

\end{document}
